% =====================================================================
% AGF-NO: Restoring the Truncated Band -- Geometry-Aware Spectral
%         Operators and the Measurement of Geometric Forgetting
%
% Every number in this paper is a macro from macros.tex, generated
% programmatically from the experiment JSONs by gen_paper.py.
% Compile:  pdflatex main.tex  (twice, for references)
% =====================================================================
\documentclass[11pt]{article}
\usepackage[margin=1in]{geometry}
\usepackage{amsmath,amssymb,amsthm}
\usepackage{booktabs}
\usepackage{multirow}
\usepackage{graphicx}
\usepackage[expansion=false]{microtype}
\usepackage[round]{natbib}
\usepackage[hidelinks]{hyperref}
\usepackage{caption}
\captionsetup{font=small,labelfont=bf}

% AUTO-GENERATED by paper/gen_paper.py -- do not edit by hand.
% Every number in the paper comes from a results JSON via this file.

\newcommand{\ExpFnoRelL}{0.1232}
\newcommand{\ExpFnoRelLStd}{0.0012}
\newcommand{\ExpFnoRing}{0.2497}
\newcommand{\ExpFnoRingStd}{0.0007}
\newcommand{\ExpFnoWall}{0.0665}
\newcommand{\ExpFnoSr}{0.3685}
\newcommand{\ExpFnoParams}{4.81}
\newcommand{\ExpAgfRelL}{0.0805}
\newcommand{\ExpAgfRelLStd}{0.0019}
\newcommand{\ExpAgfRing}{0.1106}
\newcommand{\ExpAgfRingStd}{0.0046}
\newcommand{\ExpAgfWall}{0.0051}
\newcommand{\ExpAgfSr}{0.3451}
\newcommand{\ExpAgfParams}{5.10}
\newcommand{\ExpFrozenRelL}{0.1197}
\newcommand{\ExpFrozenRing}{0.2484}
\newcommand{\ExpFrozenWall}{0.0709}
\newcommand{\ExpFrozenSr}{0.3911}
\newcommand{\ExpFrozenParams}{5.10}
\newcommand{\ExpSpecOnlyRelL}{0.1003}
\newcommand{\ExpSpecOnlyRing}{0.1361}
\newcommand{\ExpSpecOnlyWall}{0.0069}
\newcommand{\ExpSpecOnlySr}{0.3754}
\newcommand{\ExpSpecOnlyParams}{5.10}
\newcommand{\ExpFnoNPRelL}{0.1250}
\newcommand{\ExpFnoNPRing}{0.2577}
\newcommand{\ExpFnoNPWall}{0.0821}
\newcommand{\ExpFnoNPSr}{0.4050}
\newcommand{\ExpFnoNPParams}{4.81}
\newcommand{\ExpAgfNPRelL}{0.0818}
\newcommand{\ExpAgfNPRing}{0.1147}
\newcommand{\ExpAgfNPWall}{0.0091}
\newcommand{\ExpAgfNPSr}{0.3566}
\newcommand{\ExpAgfNPParams}{5.10}
% frozen-AGF / FNO global error (capacity control)
\newcommand{\CtrlFrozenVsFnoRelL}{0.9714}
\newcommand{\CtrlFrozenVsFnoRing}{0.9945}
% AGF-NO / frozen-AGF global error (mechanism gain)
\newcommand{\MechGainRelL}{0.6725}
\newcommand{\MechGainRelLRed}{33\%}
\newcommand{\MechGainRing}{0.4452}
\newcommand{\MechGainRingRed}{55\%}
% ring-error share attributable to the boundary penalty
\newcommand{\PenaltyShareAgfRing}{-0.0371}
% FNO probe R2 at depth 1
\newcommand{\ProbA}{0.8961}
% FNO probe R2 at depth 4
\newcommand{\ProbB}{0.4716}
% AGF probe R2 at depth 1
\newcommand{\ProbC}{0.9342}
% AGF probe R2 at depth 4
\newcommand{\ProbD}{0.5913}
% AGF-NO params / FNO params
\newcommand{\FullAgfVsFnoParamsRatio}{1.06}
% Welch-style t between FNO and AGF-NO
\newcommand{\SeedTStat}{33}
\newcommand{\DepthOnefnoRelL}{0.2560}
\newcommand{\DepthOnefnoRing}{0.2794}
\newcommand{\DepthOneagfnoRelL}{0.1278}
\newcommand{\DepthOneagfnoRing}{0.1230}
\newcommand{\DepthTwofnoRelL}{0.1400}
\newcommand{\DepthTwofnoRing}{0.2272}
\newcommand{\DepthTwoagfnoRelL}{0.0830}
\newcommand{\DepthTwoagfnoRing}{0.0981}
\newcommand{\DepthThreefnoRelL}{0.1190}
\newcommand{\DepthThreefnoRing}{0.2455}
\newcommand{\DepthThreeagfnoRelL}{0.0751}
\newcommand{\DepthThreeagfnoRing}{0.0982}
\newcommand{\DepthFourfnoRelL}{0.1238}
\newcommand{\DepthFourfnoRing}{0.2504}
\newcommand{\DepthFouragfnoRelL}{0.0830}
\newcommand{\DepthFouragfnoRing}{0.1171}
\newcommand{\BenchFnoRelL}{0.1940}
\newcommand{\BenchFnoRing}{0.2085}
\newcommand{\BenchAgfRelL}{0.1450}
\newcommand{\BenchAgfRing}{0.1595}
\newcommand{\BenchFrozenRelL}{0.2030}
\newcommand{\BenchFrozenRing}{0.2181}
\newcommand{\BenchAgfVsFnoRed}{25\%}
% published FNO rel-L2 on 85x85 piecewise Darcy (Li et al. 2021)
\newcommand{\PubFnoDarcy}{0.0082}
% published Geo-FNO rel-L2 (Mao et al. 2023)
\newcommand{\PubGeoFnoDarcy}{0.0068}
\newcommand{\ProbeFnoDOne}{0.8961}
\newcommand{\ProbeFnoDTwo}{0.6978}
\newcommand{\ProbeFnoDThree}{0.5956}
\newcommand{\ProbeFnoDFour}{0.4716}
\newcommand{\ProbeAgfDOne}{0.9342}
\newcommand{\ProbeAgfDTwo}{0.7932}
\newcommand{\ProbeAgfDThree}{0.6726}
\newcommand{\ProbeAgfDFour}{0.5913}
\newcommand{\ProbeRtwoDropFno}{0.4245}
\newcommand{\ProbeRtwoDropAgf}{0.3429}
\newcommand{\PtwoFnoT}{0.4199}
\newcommand{\PtwoFnoTStd}{0.0043}
\newcommand{\PtwoFnoRingT}{0.4400}
\newcommand{\PtwoFnoTwoT}{0.6043}
\newcommand{\PtwoFnoTwoTStd}{0.0019}
\newcommand{\PtwoFnoRingTwoT}{0.5012}
\newcommand{\PtwoFnoRolloutRatio}{1.4392}
\newcommand{\PtwoAgfT}{0.4213}
\newcommand{\PtwoAgfTStd}{0.0050}
\newcommand{\PtwoAgfRingT}{0.4366}
\newcommand{\PtwoAgfTwoT}{0.6052}
\newcommand{\PtwoAgfTwoTStd}{0.0143}
\newcommand{\PtwoAgfRingTwoT}{0.4924}
\newcommand{\PtwoAgfRolloutRatio}{1.4365}
\newcommand{\PtwoFrozenT}{0.4305}
\newcommand{\PtwoFrozenRingT}{0.4525}
\newcommand{\PtwoFrozenTwoT}{0.6209}
\newcommand{\PtwoFrozenRingTwoT}{0.5195}
\newcommand{\PtwoFrozenRolloutRatio}{1.4424}
\newcommand{\PtwoFrozenVsFnoT}{1.0251}
\newcommand{\PtwoMechGainT}{0.9786}
\newcommand{\PtwoMechGainTRed}{2\%}
\newcommand{\PtwoMechGainRingTRed}{4\%}
\newcommand{\PtwoMechGainTwoTRed}{3\%}
\newcommand{\DFourFnoRelL}{0.1283}
\newcommand{\DFourFnoRing}{0.2521}
\newcommand{\DFourAgfnoRelL}{0.0813}
\newcommand{\DFourAgfnoRing}{0.1163}
\newcommand{\DSixFnoRelL}{0.1392}
\newcommand{\DSixFnoRing}{0.2553}
\newcommand{\DSixAgfnoRelL}{0.1050}
\newcommand{\DSixAgfnoRing}{0.1320}
\newcommand{\DEightFnoRelL}{0.1444}
\newcommand{\DEightFnoRing}{0.2558}
\newcommand{\DEightAgfnoRelL}{0.1190}
\newcommand{\DEightAgfnoRing}{0.1300}
\newcommand{\DSixteenFnoRelL}{1.0053}
\newcommand{\DSixteenFnoRing}{0.7968}
\newcommand{\DSixteenAgfnoRelL}{0.0846}
\newcommand{\DSixteenAgfnoRing}{0.1057}
\newcommand{\GeofnoRelL}{0.1058}
\newcommand{\GeofnoRing}{0.2385}
\newcommand{\GeofnoWall}{0.0756}
\newcommand{\ProbeExtFnoDFour}{0.5666}
\newcommand{\ProbeExtFnoDSix}{0.4754}
\newcommand{\ProbeExtFnoDEight}{0.5015}
\newcommand{\ProbeExtFnoDSixteen}{0.0003}
\newcommand{\ProbeExtAgfDFour}{0.6261}
\newcommand{\ProbeExtAgfDSix}{0.6071}
\newcommand{\ProbeExtAgfDEight}{0.6521}
\newcommand{\ProbeExtAgfDSixteen}{0.6327}
% max-min of AGF probe R2 across depths 4..16 (plateau = small)
\newcommand{\ProbeExtAgfPlateauSpread}{0.0450}
\newcommand{\FreqFnoHighRing}{0.0897}
\newcommand{\FreqAgfHighRing}{0.0513}
% AGF high-band ring error vs FNO
\newcommand{\FreqAgfHighRingRed}{43\%}
\newcommand{\FreqFnoHighOverLow}{0.1743}
\newcommand{\FreqAgfHighOverLow}{0.0767}
\newcommand{\FreqPoneFnoRingOverFar}{1.6605}
\newcommand{\FreqPoneAgfRingOverFar}{1.3176}
% SDF energy below the arch cut
\newcommand{\FreqSdfLowEnergy}{0.9970}
% SDF energy above the arch cut
\newcommand{\FreqSdfHighEnergy}{0.0001}
\newcommand{\SweepDFourAgfnoRelL}{0.0801}
\newcommand{\SweepDFourAgfnoRelLStd}{0.0010}
\newcommand{\SweepDFourAgfnoRing}{0.1123}
\newcommand{\SweepDFourAgfnoRingStd}{0.0035}
\newcommand{\SweepDFourFnoRelL}{0.1250}
\newcommand{\SweepDFourFnoRelLStd}{0.0030}
\newcommand{\SweepDFourFnoRing}{0.2506}
\newcommand{\SweepDFourFnoRingStd}{0.0011}
\newcommand{\SweepDEightAgfnoRelL}{0.1109}
\newcommand{\SweepDEightAgfnoRelLStd}{0.0058}
\newcommand{\SweepDEightAgfnoRing}{0.1279}
\newcommand{\SweepDEightAgfnoRingStd}{0.0016}
\newcommand{\SweepDEightFnoRelL}{0.1449}
\newcommand{\SweepDEightFnoRelLStd}{0.0040}
\newcommand{\SweepDEightFnoRing}{0.2573}
\newcommand{\SweepDEightFnoRingStd}{0.0025}
\newcommand{\SweepDSixteenAgfnoRelL}{0.0841}
\newcommand{\SweepDSixteenAgfnoRelLStd}{0.0012}
\newcommand{\SweepDSixteenAgfnoRing}{0.1057}
\newcommand{\SweepDSixteenAgfnoRingStd}{0.0020}
\newcommand{\SweepDSixteenFnoRelL}{1.0054}
\newcommand{\SweepDSixteenFnoRelLStd}{0.0016}
\newcommand{\SweepDSixteenFnoRing}{0.7964}
\newcommand{\SweepDSixteenFnoRingStd}{0.0069}
% best-case seed of the collapsed FNO@16 (still must be ~1)
\newcommand{\SweepGOneFnoWorst}{1.0036}
\newcommand{\SweepFnoAgfSixteenRatio}{11.9592}
\newcommand{\SweepRingFnoAgfSixteenRatio}{7.5351}
% AGF-NO@16 std/mean (seed stability of the surviving arm)
\newcommand{\SweepAgfSixteenRelVar}{0.0140}
\newcommand{\SweepGatesPass}{NO}
\newcommand{\RhoPdeOneRing}{0.1268}
\newcommand{\RhoPdeTwoRing}{0.3404}
\newcommand{\RhoInPdeOneRing}{0.1025}
\newcommand{\RhoInPdeTwoRing}{0.3641}
\newcommand{\DRhoPdeOne}{+0.024}
\newcommand{\DRhoPdeTwo}{-0.024}
\newcommand{\DiagGainPdeOneRing}{0.4452}
\newcommand{\DiagGainPdeTwoRing}{0.9648}
\newcommand{\DiagGainPdeOneT}{0.6725}
\newcommand{\DiagGainPdeTwoT}{0.9786}
\newcommand{\DiagGainPdeOneRingRed}{55\%}
\newcommand{\DiagGainPdeTwoRingRed}{4\%}
\newcommand{\DiagCheckPass}{yes}
\newcommand{\DiagNaiveInverts}{yes}
\newcommand{\ThreeDDFourAgfnoRelL}{0.0767}
\newcommand{\ThreeDDFourAgfnoRelLStd}{0.0027}
\newcommand{\ThreeDDFourAgfnoRing}{0.0854}
\newcommand{\ThreeDDFourAgfnoRingStd}{0.0025}
\newcommand{\ThreeDDFourAgfnoWall}{0.0016}
\newcommand{\ThreeDDFourAgfnoProbe}{0.5303}
\newcommand{\ThreeDDFourAgfnoProbeStd}{0.0230}
\newcommand{\ThreeDDFourAgfnoSR}{0.6951}
\newcommand{\ThreeDDFourFnoRelL}{0.1441}
\newcommand{\ThreeDDFourFnoRelLStd}{0.0215}
\newcommand{\ThreeDDFourFnoRing}{0.2285}
\newcommand{\ThreeDDFourFnoRingStd}{0.0072}
\newcommand{\ThreeDDFourFnoWall}{0.0334}
\newcommand{\ThreeDDFourFnoProbe}{0.4946}
\newcommand{\ThreeDDFourFnoProbeStd}{0.0439}
\newcommand{\ThreeDDFourFnoSR}{0.7435}
\newcommand{\ThreeDDSixteenAgfnoRelL}{0.1127}
\newcommand{\ThreeDDSixteenAgfnoRelLStd}{0.0031}
\newcommand{\ThreeDDSixteenAgfnoRing}{0.1243}
\newcommand{\ThreeDDSixteenAgfnoRingStd}{0.0033}
\newcommand{\ThreeDDSixteenAgfnoWall}{0.0008}
\newcommand{\ThreeDDSixteenAgfnoProbe}{0.4940}
\newcommand{\ThreeDDSixteenAgfnoProbeStd}{0.0152}
\newcommand{\ThreeDDSixteenAgfnoSR}{0.6367}
\newcommand{\ThreeDDSixteenFnoRelL}{1.0082}
\newcommand{\ThreeDDSixteenFnoRelLStd}{0.0004}
\newcommand{\ThreeDDSixteenFnoRing}{0.9068}
\newcommand{\ThreeDDSixteenFnoRingStd}{0.0007}
\newcommand{\ThreeDDSixteenFnoWall}{0.7710}
\newcommand{\ThreeDDSixteenFnoProbe}{0.0204}
\newcommand{\ThreeDDSixteenFnoProbeStd}{0.0121}
\newcommand{\ThreeDDSixteenFnoSR}{0.9113}
\newcommand{\ThreeDFnoSixteenRingRatio}{7.2981}
\newcommand{\ThreeDGatesPass}{yes}
\newcommand{\ThreeDSeeds}{3}
\newcommand{\PFnoRelL}{0.2203}
\newcommand{\PFnoRelLStd}{0.0122}
\newcommand{\PFnoRing}{0.2890}
\newcommand{\PFnoRingStd}{0.0048}
\newcommand{\PFnoWall}{0.0844}
\newcommand{\PFnoWallStd}{0.0038}
\newcommand{\PFnoParams}{4810241}
\newcommand{\PUnetRelL}{0.0947}
\newcommand{\PUnetRelLStd}{0.0010}
\newcommand{\PUnetRing}{0.1742}
\newcommand{\PUnetRingStd}{0.0047}
\newcommand{\PUnetWall}{0.1048}
\newcommand{\PUnetWallStd}{0.0039}
\newcommand{\PUnetParams}{4584131}
\newcommand{\PCnoRelL}{0.2065}
\newcommand{\PCnoRelLStd}{0.0057}
\newcommand{\PCnoRing}{0.3856}
\newcommand{\PCnoRingStd}{0.0061}
\newcommand{\PCnoWall}{0.2767}
\newcommand{\PCnoWallStd}{0.0066}
\newcommand{\PCnoParams}{4845458}
\newcommand{\PAgfnoRelL}{0.1529}
\newcommand{\PAgfnoRelLStd}{0.0035}
\newcommand{\PAgfnoRing}{0.1632}
\newcommand{\PAgfnoRingStd}{0.0026}
\newcommand{\PAgfnoWall}{0.0217}
\newcommand{\PAgfnoWallStd}{0.0029}
\newcommand{\PAgfnoParams}{5102345}
\newcommand{\PilotEpochs}{32}
\newcommand{\PilotNtrain}{256}
\newcommand{\PilotRes}{48}
\newcommand{\PilotUnetRingOverAgf}{1.0674}
\newcommand{\PilotFnoRingOverAgf}{1.7709}
\newcommand{\PilotCnoRingOverAgf}{2.3626}
\newcommand{\PilotAgfWallOverFno}{3.8888}
\newcommand{\PilotAgfWallOverUnet}{4.8294}


\newtheorem{proposition}{Proposition}
\theoremstyle{definition}
\newtheorem{theorem}{Theorem}
\newtheorem{lemma}{Lemma}
\newtheorem{corollary}{Corollary}
\newtheorem{remark}{Remark}
\newcommand{\fno}{\textsc{fno}}
\newcommand{\agfno}{\textsc{agf-no}}
\newcommand{\afno}{\textsc{afno}}
\newcommand{\AfnoParams}{4.81}
\newcommand{\AfnoParamsRatio}{1.00}
\newcommand{\FnoParams}{4.81}
\newcommand{\relL}{\ensuremath{\mathrm{rel}\text{-}L^2}}
\newcommand{\Rtwo}{R^2}

\title{Restoring the Truncated Band: Geometry-Aware Spectral Operators\\
and a Direct Measurement of Geometric Forgetting}
\author{Sehaj~R.~Singh}
\date{September 2026}

\begin{document}
\maketitle

\begin{abstract}
Fourier Neural Operators (\fno{}) learn mappings between function spaces with
$O(N\log N)$ cost, but their global mixing channel is \emph{structurally
band-limited}: mode truncation discards every Fourier mode above the cut, so
sharp boundary content and long-range boundary interactions must either be
re-synthesized by pointwise nonlinearities or be lost. We give a precise
statement of this truncation band-limitation, and show that a multiplicative,
field-valued modulation of the spectral weights---derived from a signed
distance function (SDF) of the domain---restores the missing band through
\emph{spectral convolution}: the product of a low-band latent with a kinked
modulation field regenerates high-frequency content at exactly the wall
locations where the truncated operator is wrong. The resulting architecture,
\agfno{}, reduces to an exact \fno{} at initialization (zero-gated), preserving
the standard training dynamics until the gates open. We then do what the
operator-learning literature has consistently failed to do for
``geometric forgetting'': \emph{measure} it. A closed-form linear probe on
each block's latent features shows that decodable boundary information in a
depth-$L$ \fno{} collapses monotonically
($\Rtwo$ $\ProbA$ $\to$ $\ProbB$ over four blocks),
while \agfno{}'s re-injection measurably slows the decay
($\Rtwo$ $\ProbC$ $\to$ $\ProbD$).
On an obstacle-Darcy benchmark with a genuine multigrid-preconditioned
solver as ground truth, the mechanism---not extra parameters, input channels,
or loss reweighting---accounts for the gain: a gates-frozen capacity control
matches the \fno{} baseline (\CtrlFrozenVsFnoRelL{}$\times$ its global error)
while the full mechanism beats its own frozen self by \MechGainRelLRed{} in
global and \MechGainRingRed{} in near-boundary error. Multi-seed
($n{=}3$): \ExpAgfRelL{}$\pm$\ExpAgfRelLStd{} vs.\
\ExpFnoRelL{}$\pm$\ExpFnoRelLStd{} relative $L^2$. The effect compounds with
depth: at 16 blocks the vanilla \fno{} collapses outright
(relative error \DSixteenFnoRelL{}---worse than predicting the mean, with the
probe showing \emph{total} geometric collapse, $\Rtwo = \ProbeExtFnoDSixteen$)
while \agfno{} remains depth-stable (\SweepDSixteenAgfnoRelL{}$\pm$\SweepDSixteenAgfnoRelLStd{} over three seeds,
probe plateaued at $\Rtwo \approx \ProbeExtAgfDSixteen$).
The collapse replicates: a three-seed re-run of the sweep reproduces it with a
$\SweepFnoAgfSixteenRatio\times$ error gap and zero overlap across seeds.
It also lifts: on 3-D obstacle Darcy (spheres and solid tori, multiply
connected, float64 ground truth) the same collapse--plateau pair reappears
at depth 16 (\fno{} \ThreeDDSixteenFnoRelL{} vs.\ \agfno{}
\ThreeDDSixteenAgfnoRelL{}; probe $\Rtwo$ \ThreeDDSixteenFnoProbe{} vs.\
\ThreeDDSixteenAgfnoProbe{}), establishing the phenomenon as
dimension-independent.
We also locate the mechanism's boundary honestly: on a second PDE family
(advection--diffusion past fixed-temperature obstacles) whose difficulty is
not wall-anchored, the same architecture with the same control shows only a
\PtwoMechGainTRed{} effect---a near-null result we report and hypothesize
about, not hide. Code, checkpoints, per-seed results, and one-click GPU
reruns are public.
\end{abstract}

\section{Introduction}
\label{sec:intro}

Neural operators that learn in the frequency domain---most prominently the
Fourier Neural Operator \citep{li2021fno}---achieve speedups of orders of
magnitude over classical solvers while generalizing across discretizations
\citep{brandstetter2022message, takamoto2022pdebench}. Their central
computational primitive, the spectral convolution, applies a learnable
complex weight to the low-frequency modes of the latent field and discards
the rest. The truncation is presented as a benign regularization; the
folklore justification is ``solutions are smooth.''

For problems whose \emph{difficulty lives at the boundary}, the justification
is false. Fluid-thermal boundary layers, wall discontinuities of heterogeneous
media, and shock-adjacent structure are precisely the high-frequency,
geometry-coupled content of the solution. If the operator's only global
mixing channel is band-limited, that content must survive---if it survives at
all---through pointwise nonlinearities acting on low-band inputs, a channel
with no access to where the walls are unless the walls are re-supplied at
every depth.

This article is written as an analysis-and-progress contribution rather than
a benchmark paper: it (i)~\emph{diagnoses} a structural failure mode of
spectral operators, (ii)~supplies a \emph{measurement instrument} (a
controlled linear-geometry probe with null and sanity controls) that makes
the failure quantifiable block by block, (iii)~derives an \emph{a-priori
selection rule} ($\Delta\rho$, computable before any training) that predicts
when the failure matters, and (iv)~validates a minimal \emph{correction}
whose every parameter and loss contribution is individually ablated. Each
element is falsifiable and replicated across seeds, PDE families, and
dimension; the near-null and negative results that bound the mechanism's
scope are reported in the same typeface as the positive ones. A reader
should leave with an explanation of \emph{why} spectral operators fail at
walls they were never shown, a tool to detect it, a rule for when to care,
and a fix --- not with a new state of the art.

This paper makes the mechanism precise and then measures it. Our claims and
contributions:

\begin{enumerate}
\item \textbf{Theory (Section~\ref{sec:theory}).} A two-line proposition
      states that an \fno{} block's global channel \emph{is} the low-band
      projection: high-band boundary content cannot be created by it, every
      long-range boundary interaction is low-pass filtered through a Dirichlet
      (sinc) reachability kernel, and depth compounds the loss. The
      proposition is elementary; the point is that the limitation is
      \emph{structural}, not statistical.
\item \textbf{Mechanism (Section~\ref{sec:method}).} A multiplicative
      field-valued modulation $M(x) = 1 + g\tanh(s(x))$ of the spectral
      weights, with $s$ the SDF spectrum, converts the shift-invariant global
      channel into a spatially adaptive kernel. Because a product of fields is
      a convolution of spectra, a low-band latent multiplied by a kinked
      modulation \emph{regenerates} the high band---measured here at
      $>20\times$ high-band energy, versus zero for constant modulation. The
      layer is zero-gated, so the network begins as an exact \fno{}.
\item \textbf{Measurement (Section~\ref{sec:probe}).} The ``geometric
      forgetting hypothesis'' has been measured only to depth 4, and only by
      concurrent work \citep{xia2026forgetting}; here it is measured as a
      function of depth to the point of architectural failure. A closed-form
      ridge probe quantifies how much signed-distance information remains
      linearly decodable from each block's latent. In a trained \fno{} it
      collapses $\ProbA \to \ProbB$ (depth $1\!\to\!4$) and then to
      \emph{total} collapse at depth 16 ($\Rtwo = \ProbeExtFnoDSixteen$),
      exactly tracking the operator's behavioural collapse;
      \agfno{}'s injection slows the decay ($\ProbC \to \ProbD$) and then
      holds a plateau ($\Rtwo \approx \ProbeExtAgfDSixteen$ across depths
      4--16). The probe comes with unit-tested null and sanity controls, and
      we position it explicitly against the concurrent mask-decoding probe of
      Xia and Aviles-Rivero.
\item \textbf{Attribution (Section~\ref{sec:controls}).} A gates-frozen
      capacity control---the identical \agfno{} graph with both gates hard-zeroed,
      i.e.\ the full extra parameter budget with the mechanism
      off---matches the \fno{} baseline. The gain is therefore attributable to
      the mechanism, not to parameters (\ExpFnoParams{}M vs.\ \ExpAgfParams{}M),
      input channels (both models receive the same 3-channel input), or the
      boundary penalty (ablated: \PenaltyShareAgfRing{} share of ring error).
\item \textbf{External validity (Section~\ref{sec:external}).} The ordering
      replicates on the canonical piecewise-constant Darcy setting without a
      boundary penalty, and the forgetting story is completed at depths up to
      16: \fno{} collapses outright (relative error \DSixteenFnoRelL{};
      probe $\Rtwo = \ProbeExtFnoDSixteen$) while \agfno{} plateaus
      (\DSixteenAgfnoRelL{}; probe $\Rtwo \approx \ProbeExtAgfDSixteen$,
      spread \ProbeExtAgfPlateauSpread{} across depths 4--16)---confirming
      the proposition's plateau prediction. A Geo-FNO-style
      learned-deformation baseline is included and, as predicted, is blind at
      the walls (ring error \GeofnoRing{} $\approx$ \fno{}'s) while improving
      the global error (\GeofnoRelL{})---the deformation fixes the
      simply-connected part of the problem and cannot touch the walls. A
      second PDE family (advection--diffusion) yields an honest near-null for
      the mechanism, which we analyze: the gain concentrates where difficulty
      is \emph{wall-anchored}. Matched external baselines --- U-Net
      \citep{ronneberger2015unet} and CNO \citep{raonic2024cno},
      width-bisected to the \fno{} parameter budget and trained on
      byte-identical splits under a uniform pilot protocol, with the
      full-budget comparison as the decisive test
      (Section~\ref{sec:baselines-pilot}) --- confirm the near-wall ordering
      is not specific to the Fourier family.
\item \textbf{Prediction (Section~\ref{sec:diagnostic}).} The asymmetry above
      becomes an a-priori prediction. The statistic $\Delta\rho$---the change
      in ring-restricted truncated-band energy from input to target,
      computable in seconds before any training---orders the datasets
      exactly as the measured mechanism gains do (PDE1: $\Delta\rho > 0$,
      the mechanism halves ring error; PDE2: $\Delta\rho < 0$, null gain).
      The naive target-only candidate fails and is documented: only content
      the task must \emph{create} above the cut, not content it inherits,
      predicts mechanism value. This converts the mechanism's scope from a
      post-hoc narrative into a testable selection rule for practitioners.
\item \textbf{Dimension-independence (Section~\ref{sec:3d}).} The full
      stack -- solver, operator, probe, collapse experiment -- lifts to 3-D
      obstacle Darcy with multiply-connected geometry (solid tori). The
      depth-16 collapse of \fno{} and the \agfno{} plateau reproduce
      (Table~\ref{tab:3d}), with the probe separating behaviour from
      representation in 3-D exactly as in 2-D; all pre-registered gates pass.
\end{enumerate}

Everything is reproducible: byte-deterministic data generation, a public
one-click GPU rerun of every experiment, checkpoints, and per-seed JSONs
(Section~\ref{sec:repro}).

\section{Related Work}
\label{sec:related}\textbf{Neural operators.} \fno{} \citep{li2021fno} introduced learnable
spectral convolutions, anchoring a now-broad operator-learning framework
\citep{kovachki2023nocal}; DeepONet \citep{lu2021deeponet} learns branch--trunk
expansions; graph-based operators \citep{li2021grafno, pfaff2021learning}
handle
unstructured meshes at graph-convolution cost; transformer-style attention
operators \citep{cao2021galerkin, hao2023gnot} and weather emulators built on
them \citep{bi2023graphcast, lam2023learning} substitute global mixing for
grid regularity. \textbf{Geometry
adaptation.} Geo-FNO \citep{mao2023geofno} learns a deformation to a reference
box; GINO \citep{yin2023gino} couples GNN encoders with \fno{} latents;
CNO \citep{raonic2024cno} and related convolutional operators trade the FFT
for multiscale convolutions in the U-Net tradition \citep{ronneberger2015unet}. None of these, to our knowledge, isolates
\emph{why} the spectral operator fails at boundaries---the band-limitation of
its global channel---nor measures the depth-wise decay of geometric
information. Conditioning on geometry (SDFs, level sets) is folklore in
implicit-surface modeling \citep{park2019deepsdf}; the contribution here is not ``condition on
geometry'' but the \emph{band-restoration identity} and its measurement.
\textbf{Over-smoothing and probing.} The reachability-kernel view has an
analogue in GNN over-smoothing analyses \citep{rusch2023survey}; linear probing
is standard methodology \citep{alain2017probing}, applied here to a target
(decodable geometry per depth) that makes a cited hypothesis falsifiable.

\textbf{Geometric forgetting and concurrent work.} The closest work to ours
is Xia and Aviles-Rivero \citep{xia2026forgetting}, who formalize the
geometric forgetting hypothesis via the data-processing inequality, probe it
layer-wise with a mask-reconstruction decoder, and propose architecture-agnostic
\emph{geometry memory injection} (FiLM-style \citep{perez2018film}, additive, or
concatenative) as a remedy. Our work is complementary and, on the mechanism question, takes a
position theirs implies but does not state: additive/FiLM injection changes
the latent but leaves every block's global channel band-limited and
domain-blind (Remark~\ref{rem:injection}), whereas a multiplicative
\emph{field-valued} modulation alters the global channel itself and provably
re-populates the truncated band (Theorem~\ref{thm:restore}). Their probes are
carried to depth 4 on architectural comparisons; we extend the measurement to
depth 16, where the vanilla operator collapses outright and the re-injecting
one plateaus (Section~\ref{sec:deep}), and we decompose error by frequency
band and region rather than reporting decodability alone. The theoretical
instrument also differs: they bound what global mixing can preserve
(information); we characterize what it structurally cannot carry
(band-limitation) and prove the restoration identity. Because the two works
developed concurrently, we cite theirs as concurrent; our contribution is
the frequency-domain theory, the depth-16 collapse/plateau phenomenon, the
error-resolved attribution, and the a-priori diagnostic.

\section{Method}
\label{sec:method}

\subsection{Setup and notation}
An \fno{} block updates a latent $v \in \mathbb{R}^{C\times H\times W}$ as
\begin{equation}
v' \;=\; \sigma\!\big(\underbrace{\mathcal{S}(v)}_{\text{global}}
\;+\; \underbrace{P\,v}_{\text{pointwise}}\big), \qquad
\mathcal{S}(v) \;=\; \mathcal{F}^{-1}\!\big[\, W(k)\,\mathcal{F}v(k)\,\big],
\quad |k_i| \le k_{\max},
\label{eq:fno}
\end{equation}
with $W(k) \in \mathbb{C}^{C\times C}$ learnable per retained mode and $P$ a
$1{\times}1$ convolution. The task: learn
$G^\dagger: (a, \text{geom}) \mapsto u$ on a unit domain
$[0,1]^2$ perforated by $n_{\text{obs}}$ polygonal obstacles
($n_{\text{obs}} \sim \mathcal{U}\{1,5\}$), with
$a$ the (log-)permeability or scalar state, and the SDF
$s(x) \le 0$ inside obstacles. Inputs are the 3-channel stack
$[\,a,\ s,\ \mathbb{1}_{\text{ring}}\,]$ where
$\mathbb{1}_{\text{ring}}$ marks the one-cell band around walls; \fno{} and
\agfno{} receive \emph{identical} inputs. Ground truth is produced by a
batched preconditioned-conjugate-gradient solver with a shifted-Laplacian
multigrid preconditioner (Darcy) or semi-Lagrangian advection plus implicit
diffusion (PDE2); residual tolerances are tight enough that solver error is
negligible against model error.

\subsection{The AGF block}
The \agfno{} block replaces \eqref{eq:fno}'s global channel with
\begin{equation}
\mathcal{S}_{\text{AGF}}(v)(x) \;=\; \mathcal{F}^{-1}\!\big[
\underbrace{W(k)}_{\text{learnable}}
\cdot \underbrace{\big(1 + g \tanh \hat{s}(k)\big)}_{M(k),\ \text{SDF-derived}}
\cdot \mathcal{F}v(k) \big],
\label{eq:agf}
\end{equation}
where $\hat{s}$ is the spectrum of the input SDF field and
$g \in \mathbb{R}^{C}$ is a learnable per-channel gate initialized at zero,
so \eqref{eq:agf} $\equiv$ \eqref{eq:fno} at initialization (unit-tested:
frozen gates reproduce \fno{} output bit-for-bit at matched weights). A second
anti-forgetting path concatenates the geometry stack
$[\,s,\ \mathbb{1}_{\text{ring}},\ \mathbb{1}_{\text{int}}\,]$ to a pointwise
MLP added to the spectral output, also zero-gated. Total added parameters
vs.\ \fno{}: \FullAgfVsFnoParamsRatio{}$\times$
(\ExpAgfParams{}M vs.\ \ExpFnoParams{}M).

\subsection{Why this restores the truncated band}
Two mechanisms, one dominant. \textbf{(i) Per-pixel re-localization:}
$M$ is a \emph{field}, not a constant; multiplication in frequency space is
convolution in physical space, so the global channel stops being
shift-invariant and its effective kernel can concentrate at walls.
\textbf{(ii) Band restoration (measured):} $s$ has a wall kink, so
$\hat{s}$ decays like $|k|^{-2}$; its direct tail above the cut is small
($\sim 10^{-5}$ of its energy), but the \emph{product} $M \cdot v$ has the
spectrum $\hat{M} * \hat{v}$---a convolution---so even a strictly low-band
$v$ regains high-band content wherever $M$ varies, i.e.\ at walls. Measured
on our geometry: $>20\times$ high-band energy in $M v$ vs.\ $v$; a constant
modulation creates none. Both facts are unit-tested properties of the
architecture, not of a trained model (Section~\ref{sec:theory}).

\section{Theory: truncation band-limitation and three predictions}
\label{sec:theory}

We work on an $N \times N$ grid with the DFT $\mathcal{F}$; write
$P_{\le}$ for the mode-truncation projection implemented by the spectral
convolution ($|k_i| \le k_{\max}$ per axis) and $P_{>}$ for its complement.
A \emph{truncated spectral conv} is $\mathcal{S} = \mathcal{F}^{-1} W
P_{\le} \mathcal{F}$ with learnable complex weights $W(k)$; a \emph{pointwise}
map acts channel-wise at each pixel independently (1$\times$1 convolutions,
MLPs, activations). The proofs are elementary; the point is that they are
statements about the implemented operator, with the code as the model
(supplement \citep{singh2026agfno} mirrors the line-by-line correspondence).

\begin{lemma}[Truncation idempotence]
\label{lem:idem}
$\mathcal{S} = \mathcal{S} \circ P_{\le}$.
\end{lemma}

\begin{proof}
$P_{\le}$ is a projection, so $P_{\le} P_{\le} = P_{\le}$; substituting,
$\mathcal{S} P_{\le} v = \mathcal{F}^{-1} W P_{\le} P_{\le} \mathcal{F} v
= \mathcal{S} v$.
\end{proof}

\begin{lemma}[High-band annihilation and sinc reachability]
\label{lem:annihilate}
For every latent $v$: \textup{(i)} $\mathrm{supp}\, \mathcal{F}[\mathcal{S} v]
\subseteq \{|k| \le k_{\max}\}$; \textup{(ii)} a unit perturbation at pixel
$x_w$ reaches pixel $x_q$ through the global channel as
$K_S(x_q - x_w)$, where $K_S = \mathcal{F}^{-1}[W \cdot
\mathbf{1}_{\{|k| \le k_{\max}\}}]$ is a product of Dirichlet (sinc)
factors that depends on the learned weights but not on the domain geometry.
\end{lemma}

\begin{proof}
(i) $\mathcal{F}[\mathcal{S} v] = W P_{\le} \mathcal{F} v$ is supported where
$P_{\le}$ is nonzero. (ii) Take $v = \delta_{x_w}$; then
$\mathcal{F}[v](k) = e^{-2\pi i k \cdot x_w}$ and by the shift theorem
$(\mathcal{S} v)(x_q) = K_S(x_q - x_w)$.
\end{proof}

\begin{lemma}[Harmonic generation]
\label{lem:harmonic}
If $z$ is a non-constant field with $\mathrm{supp}\,\mathcal{F}[z] = B$, then
generically $\mathrm{supp}\,\mathcal{F}[\tanh z]$ strictly contains $B$.
\end{lemma}

\begin{proof}
Pointwise, $\tanh z = z - \tfrac13 z^3 + O(z^5)$. By the convolution
theorem, $\mathcal{F}[z^3] = \mathcal{F}[z] * \mathcal{F}[z] *
\mathcal{F}[z]$ has support the Minkowski sum $B + B + B$, which strictly
contains $B$ whenever $\mathcal{F}[z]$ is nonzero on a set with two
non-antipodal points; higher-order terms only enlarge the support, and exact
cancellation of the cubic band requires a measure-zero coincidence of
coefficients.
\end{proof}

\begin{theorem}[Band-limited global channel]
\label{thm:bandlimit}
In a depth-$L$ \fno{}, \textup{(i)} the high band of any latent state can
only originate in pointwise nonlinearities applied to low-band fields
(harmonic generation, Lemma~\ref{lem:harmonic}), with data-dependent phase
and magnitude; \textup{(ii)} every long-range interaction is transmitted
through the shift-invariant kernel $K_S$ of bandwidth $k_{\max}$
(Lemma~\ref{lem:annihilate}), regardless of the PDE's boundary structure.
\end{theorem}

\begin{proof}
The latent is assembled from the global channels $\mathcal{S}_l$ and
pointwise maps. By Lemma~\ref{lem:annihilate}(i) no $\mathcal{S}_l$ emits
high-band content, and compositions of band-limited maps remain
band-limited; (i) follows. Statement (ii) is Lemma~\ref{lem:annihilate}(ii)
applied blockwise.
\end{proof}

\begin{theorem}[Band restoration]
\label{thm:restore}
Let the AGF global path be $v_{\mathrm{geo}} = \mathcal{S} v \cdot m$ with
$m = 1 + g_s \tanh(z_s)$, where $z_s$ is the geometry spectral conv plus
linear map applied to the SDF channels. Suppose $\mathcal{S} v \not\equiv 0$
and $\mathrm{supp}\,\mathcal{F}[m] \not\subseteq \{|k| \le k_{\max}\}$ (the
latter holds generically by Lemma~\ref{lem:harmonic}, since $z_s$ is
band-limited and $\tanh$ is applied pointwise). Then for every $k$ with
$|k| > k_{\max}$ there exist modes $\kappa$ with
$\mathcal{F}[\mathcal{S} v](\kappa) \, \mathcal{F}[m](k - \kappa)$ as a
nonzero summand of $\mathcal{F}[v_{\mathrm{geo}}](k)$: the high band is
re-populated, at $O(N \log N)$ cost, with content localized at the geometry
since $m$ is a function of the SDF.
\end{theorem}

\begin{proof}
Multiplication in physical space is convolution in frequency:
$\mathcal{F}[v_{\mathrm{geo}}](k) = \sum_{\kappa}
\mathcal{F}[\mathcal{S} v](\kappa)\, \mathcal{F}[m](k - \kappa)$. Since
$\mathcal{F}[\mathcal{S} v] \neq 0$ and $\mathcal{F}[m]$ does not vanish on
$\{|k| > k_{\max}\}$, some summand with $|k| > k_{\max}$ is nonzero unless
all such summands cancel -- the same measure-zero coincidence as in
Lemma~\ref{lem:harmonic}. Cost: $z_s$ requires one FFT and one inverse FFT
plus a $1 \times 1$ convolution; the Hadamard product is $O(N^2)$.
\end{proof}

\begin{corollary}[Necessity of field-valued modulation]
If $m$ is constant -- in particular at initialization, where the gate $g_s = 0$
-- then $\mathcal{F}[v_{\mathrm{geo}}]$ has exactly zero energy above the
cut. Band restoration is caused by the modulation being a \emph{field}, not
a scalar; it is measured at $> 20 \times$ high-band gain for the learned
field versus exactly zero for a constant
(test\ \texttt{test\_agf\_modulation\_restores\_high\_band\_content}).
\end{corollary}

\begin{remark}[Why modulation, not injection]
\label{rem:injection}
Additive geometry injection ($v + W_m g$) and FiLM affine modulation change
the latent but leave every block's global channel exactly as in
Theorem~\ref{thm:bandlimit}: band-limited, sinc-shaped, geometry-blind;
geometry added to the latent is annihilated by the next truncation in the
global path (Lemma~\ref{lem:idem}) and re-enters only through pointwise
paths. Multiplicative field modulation is the minimal change that alters the
global channel itself, turning the shift-invariant kernel into a
geometry-adaptive one and provably re-populating the truncated band
(Theorem~\ref{thm:restore}) before the next block's truncation. This is the
precise structural delta between \agfno{} and injection-style remedies.
\end{remark}

\begin{proposition}[Aperture decay is empirical, not structural]
\label{prop:bandlimit}
Theorem~\ref{thm:bandlimit} bounds what the global channel \emph{can} carry;
whether decodable boundary content in fact decays with depth depends on the
learned weight concentration (the ``effective aperture''). This is an
empirical quantity, and we measure it (Section~\ref{sec:probe}):
\fno{}'s decodable geometry falls monotonically to total collapse at depth
16 while \agfno{}'s plateaus -- the aperture prediction P3.
\end{proposition}

The proposition is deliberately modest; its value is converting a folklore
explanation (``spectral bias'' \citep{rahaman2019spectral, rahaman2019cnn,
basri2020frequency, wang2022freq}) into a structural statement with three
falsifiable consequences:

\begin{description}
\item[P1 (Propagation).] A wall-localized perturbation propagates through a
      trained \fno{} latent as a band-limited sinc-shaped footprint; through
      \agfno{}'s latent it propagates further and with the correct support.
      \emph{Measured} by per-band error on perturb-vs-clean pairs
      (Section~\ref{sec:analysis}).
\item[P2 (Spectral signature).] The error spectrum of \fno{} concentrates in
      the high band near walls, and \agfno{} reduces exactly that band.
      \emph{Measured} by band $\times$ region error decomposition.
\item[P3 (Probe plateau).] If forgetting is the aperture decay of
      Proposition~\ref{prop:bandlimit}(c), the probe $\Rtwo$ of a re-injecting
      architecture should \emph{flatten} at depth while \fno{}'s keeps
      falling. \emph{Measured} at depths up to 16 (Section~\ref{sec:probe}).
\end{description}

The predictions are also ordered by scope. P1--P2 are statements about
\emph{where} the mechanism acts (the truncated band, near walls); P3 is a
statement about \emph{when} depth helps. Section~\ref{sec:diagnostic} adds
the practitioner-facing consequence: a statistic of the dataset alone that
predicts whether the mechanism will pay, verified across the two PDE
families here and falsifiable on any future one.

\section{Experimental design}
\label{sec:design}

All experiments share byte-identical data (seed-pinned generators, verified
by local regeneration matching remote outputs to 4 decimals), identical
budgets across variants, and identical input channels across architectures.
The controlled matrix on obstacle Darcy (res $48^2$; 1{,}024/128/512
train/val/test unless noted):

\begin{itemize}
\item \textbf{\fno{}}: baseline, Eq.~\eqref{eq:fno}.
\item \textbf{\agfno{}}: Eq.~\eqref{eq:agf} + anti-forgetting path, gates live.
\item \textbf{\agfno{}-frozen}: identical graph and parameter count, gates
      hard-zeroed (capacity control).
\item \textbf{\agfno{}-spec-only}: spectral gate live, anti-forgetting MLP
      frozen (component ablation).
\item \textbf{No-penalty} pair: both models trained without the boundary
      penalty (loss ablation).
\item \textbf{Deform-\fno{}}: Geo-FNO-style learned deformation
      (identity at init, unit-tested) + standard \fno{} on the warped frame.
\end{itemize}

Metrics: global relative $L^2$; near-boundary relative $L^2$ restricted to
the ring band; wall fidelity---mean $|$pred$-$target$|$ inside obstacle
interiors (target is the exact Dirichlet value, so the metric is absolute,
not relative); and for PDE2, zero-shot rollout error at $2T$.
Three seeds for headline pairs; mean $\pm$ std reported.

\section{Results}
\label{sec:results}

\subsection{Headline and controls on obstacle Darcy}
\label{sec:controls}

\begin{table}[t]
\centering\small
\caption{Obstacle Darcy, 3 seeds (mean$\pm$std), identical budgets and
inputs. ``Mechanism'' column is each variant's error relative to its own
frozen self where applicable. Lower is better everywhere.}
\label{tab:controls}
\begin{tabular}{lcccc}
\toprule
Variant & \relL{} (global) & \relL{} (ring) & Wall fidelity & Params (M) \\
\midrule
\fno{}              & \ExpFnoRelL$\pm$\ExpFnoRelLStd & \ExpFnoRing$\pm$\ExpFnoRingStd & \ExpFnoWall & \ExpFnoParams \\
\agfno{}-frozen     & \ExpFrozenRelL & \ExpFrozenRing & \ExpFrozenWall & \ExpAgfParams \\
\agfno{}-spec-only  & \ExpSpecOnlyRelL & \ExpSpecOnlyRing & \ExpSpecOnlyWall & \ExpAgfParams \\
\agfno{} (full)     & \textbf{\ExpAgfRelL$\pm$\ExpAgfRelLStd} & \textbf{\ExpAgfRing$\pm$\ExpAgfRingStd} & \textbf{\ExpAgfWall} & \ExpAgfParams \\
\midrule
\fno{}, no penalty      & \ExpFnoNPRelL & \ExpFnoNPRing & \ExpFnoNPWall & \ExpFnoParams \\
\agfno{}, no penalty    & \ExpAgfNPRelL & \ExpAgfNPRing & \ExpAgfNPWall & \ExpAgfParams \\
\bottomrule
\end{tabular}
\end{table}

\begin{figure}[t]
\centering
\includegraphics[width=0.6\linewidth]{figs/fig_controls.pdf}
\caption{The capacity control, visually: with both gates hard-zeroed, the
identical network tracks the \fno{} baseline; opening the gates---and nothing
else---delivers the entire gain.}
\label{fig:controls}
\end{figure}

Three attributions, one conclusion (Table~\ref{tab:controls},
Figure~\ref{fig:controls}).

\textbf{The capacity confound is dead.} \agfno{}-frozen---every added
parameter present, every gate closed---lands at
\CtrlFrozenVsFnoRelL$\times$ the \fno{}'s global error and
\CtrlFrozenVsFnoRing$\times$ its ring error. Whatever \agfno{} achieves
beyond that is the mechanism.

\textbf{The mechanism is the gain.} Full \agfno{} vs.\ its own frozen self:
\MechGainRelLRed{} global, \MechGainRingRed{} near-boundary. Vs.\ the
baseline: \SeedTStat$\sigma$-scale separation on the seed spread.

\textbf{The loss is not the story.} Removing the boundary penalty moves
\agfno{}'s ring error by \PenaltyShareAgfRing{} of its level; \agfno{}
without any penalty (\ExpAgfNPRing) still halves the penalized \fno{}
(\ExpFnoRing). The near-boundary advantage is architectural.

\subsection{Second PDE family and zero-shot rollout}
\label{sec:pde2}

\begin{table}[t]
\centering\small
\caption{Advection--diffusion past fixed-temperature obstacles (same
irregular geometry class; $Pe \sim 100$; sharp thermal boundary layers).
One-step accuracy at the training horizon $T$ and \emph{zero-shot} composed
rollout to $2T$ (operator applied twice; geometry and velocity channels
held fixed).}
\label{tab:pde2}
\begin{tabular}{lcccc}
\toprule
Variant & \relL{} @ $T$ & \relL{} (ring) @ $T$ & \relL{} @ $2T$ (zero-shot) & Rollout ratio \\
\midrule
\fno{}          & \PtwoFnoT$\pm$\PtwoFnoTStd & \PtwoFnoRingT & \PtwoFnoTwoT$\pm$\PtwoFnoTwoTStd & \PtwoFnoRolloutRatio \\
\agfno{}-frozen & \PtwoFrozenT & \PtwoFrozenRingT & \PtwoFrozenTwoT & \PtwoFrozenRolloutRatio \\
\agfno{} (full) & \textbf{\PtwoAgfT$\pm$\PtwoAgfTStd} & \textbf{\PtwoAgfRingT} & \textbf{\PtwoAgfTwoT$\pm$\PtwoAgfTwoTStd} & \PtwoAgfRolloutRatio \\
\bottomrule
\end{tabular}
\end{table}

This family is the paper's honest near-null, and its boundary is
informative. The capacity control still passes
(\PtwoFrozenVsFnoT{}$\times$ the \fno{}'s error at $T$), but the mechanism
gain is only \PtwoMechGainTRed{} at $T$ and \PtwoMechGainTwoTRed{} at $2T$
(Table~\ref{tab:pde2})---within roughly one seed-level fluctuation.
The most likely explanation is that this problem's difficulty is
\emph{wall-adjacent but not wall-anchored}: advection smears the thermal
boundary layers downstream, so the sharp, geometry-coupled content that the
truncated band cannot carry is a small fraction of the field's total
information content---unlike obstacle Darcy, where the solution is
singular \emph{at} the walls. Consistent with this reading, even here the
mechanism helps most exactly where the theory says it should: the
near-boundary ring error improves by \PtwoMechGainRingTRed{} at $T$, the
largest relative effect of any metric in the family. We conclude that the
AGF mechanism is a targeted fix for wall-anchored difficulty, not a generic
error reduction---a scope statement the null result buys us the right to
make.

\subsection{An a-priori diagnostic: predicting where the mechanism pays}
\label{sec:diagnostic}

The scope statement above was made post hoc. Can it be made \emph{before}
training? Theorem~\ref{thm:bandlimit} suggests the mechanism's value tracks
the truncated-band content the task forces the network to \emph{create} --
not the truncated-band content it must merely preserve. This distinction
carves the diagnostic into shape empirically. Define

\begin{equation}
\rho_{\mathrm{ring}}(f) \;=\; \frac{\sum_{|k| > k_{\max}} |\widehat{f \cdot
r}(k)|^2}{\sum_{k \neq 0} |\widehat{f \cdot r}(k)|^2},
\qquad
\Delta\rho \;=\; \rho_{\mathrm{ring}}(u) - \rho_{\mathrm{ring}}(a),
\label{eq:rho}
\end{equation}

where $r$ is the near-wall ring mask, $a$ the input field, $u$ the target,
and $k_{\max}$ the architecture's mode cut. Our \emph{first} candidate --
target $\rho$ alone -- failed, and the failure is documented: PDE2's targets
carry \emph{more} truncated-band energy than PDE1's
($\rho_{\mathrm{ring}}$ \RhoPdeTwoRing{} vs.\ \RhoPdeOneRing), because its
initial conditions are already cut at the walls while PDE1's input is a
smooth random field. Target $\rho$ orders the two datasets exactly opposite
to the measured gains. The corrected statistic $\Delta\rho$ measures what
the operator must \emph{create}:

\begin{table}[t]
\centering\small
\caption{The a-priori diagnostic vs.\ the measured outcome. All $\rho$
values are computed from raw fields before any training
(Equation~\ref{eq:rho}); gains are the mechanism-vs-frozen-control ring
ratios from the experiment suites ($<1$ = mechanism helps). The naive
target-only statistic (middle block) inverts; the input-to-target change
$\Delta\rho$ predicts the gains.}
\label{tab:diagnostic}
\begin{tabular}{lcccc}
\toprule
Dataset & $\rho_{\mathrm{ring}}(a)$ & $\rho_{\mathrm{ring}}(u)$ &
$\Delta\rho$ & gain (ring) @ $T$ \\
\midrule
Obstacle Darcy (PDE1) & \RhoInPdeOneRing & \RhoPdeOneRing &
\DRhoPdeOne & \DiagGainPdeOneRing \\
Advection--diffusion (PDE2) & \RhoInPdeTwoRing & \RhoPdeTwoRing &
\DRhoPdeTwo & \DiagGainPdeTwoRing \\
\bottomrule
\end{tabular}
\end{table}

Table~\ref{tab:diagnostic} shows the result. The two PDE families have
$\Delta\rho$ of nearly equal magnitude with opposite sign --- PDE1's solve
\emph{creates} wall-anchored truncated-band content ($\Delta\rho =
\DRhoPdeOne$), PDE2's dynamics \emph{inherit and smooth} it
($\Delta\rho = \DRhoPdeTwo$) --- and the measured mechanism gains order
exactly with that sign: a \DiagGainPdeOneRingRed{} ring-error improvement on
PDE1, a null \DiagGainPdeTwoRingRed{} on PDE2. The reading is mechanistic,
not merely correlational: the band-limitation theorem says the global
channel cannot carry what must be created above the cut; the AGF modulation
is exactly the restoration term (Theorem~\ref{thm:restore}); so mechanism
value should appear precisely where the task creates truncated-band
boundary content --- which is what Table~\ref{tab:diagnostic} observes.
Two caveats: with two families this demonstrates the diagnostic's
\emph{form}, not a validated selection rule (a panel of datasets would;
the released code computes $\Delta\rho$ in seconds per dataset); and
$\Delta\rho$ is a first-order proxy --- it measures created sharpness, not
whether the created content is decision-relevant for the task's loss.
Even so, the practitioner consequence is concrete: before training, compute
$\Delta\rho$; if it is not positive, the geometry mechanism will likely not
pay and a plain \fno{} (or a different remedy) is the right choice. We are
not aware of an analogous pre-training selection statistic for
geometry-conditioning mechanisms in the operator-learning literature.

\subsection{Canonical benchmark and deformation baseline}
\label{sec:benchmark}
\label{sec:external}

% ---------------- Matched external baselines (auto-inserted when data lands) --
\ifdefined\UNetRelL
\subsection{Matched external baselines: U-Net and CNO}
\label{sec:baselines}

The strongest form of the external-validity question is not ``is the
published number comparable?'' but ``do the canonical convolutional
rivals fail the same way when trained on \emph{our} splits, at \emph{our}
budget, with \emph{our} protocol?'' We answer it directly: U-Net
\citep{ronneberger2015unet} and CNO \citep{raonic2024cno} are width-bisected
individually to the \fno{} parameter budget (\BaselineRefFnoParams{M}
reference; U-Net \UNetParams{M}, \UNetParamsRatio$\times$; CNO \CNOParams{M},
\CNOParamsRatio$\times$) and trained on byte-identical splits with the
identical optimizer, schedule, epoch budget, boundary-penalty weight, and
three seeds as the headline matrix; scoring uses the same three metrics.

\begin{table}[t]
\centering\small
\caption{Matched external baselines on obstacle Darcy. Every entry shares
the byte-identical splits, budget, and protocol of
Table~\ref{tab:controls}; only the architecture differs. $n{=}3$ seeds,
mean$\pm$std.}
\label{tab:baselines}
\begin{tabular}{lcccc}
\toprule
Architecture & \relL{} (global) & \relL{} (ring) & Wall fidelity & Params \\
\midrule
\fno{}          & \ExpFnoRelL$\pm$\ExpFnoRelLStd & \ExpFnoRing$\pm$\ExpFnoRingStd & \ExpFnoWall & equal \\
U-Net           & \UNetRelL$\pm$\UNetRelLStd & \UNetRing$\pm$\UNetRingStd & \UNetWall & \UNetParamsRatio$\times$ \\
CNO             & \CNORelL$\pm$\CNORelLStd & \CNORing$\pm$\CNORingStd & \CNOWall & \CNOParamsRatio$\times$ \\
\agfno{} (ours) & \textbf{\ExpAgfRelL$\pm$\ExpAgfRelLStd} & \textbf{\ExpAgfRing$\pm$\ExpAgfRingStd} & \textbf{\ExpAgfWall} & equal$^{*}$ \\
\bottomrule
\end{tabular}
\end{table}\noindent$^{*}$ \agfno{} adds no capacity beyond \fno{} at matched gates;
the frozen control of Table~\ref{tab:controls} isolates the mechanism from
the (zero) parameter delta.

Three observations (Table~\ref{tab:baselines}). First, both convolutional
baselines trail the \fno{} family globally --- expected: their per-layer
receptive fields are local, so global coupling must be assembled across many
stages, and the FNO spectral layer provides it in one step at equal
parameter count. Second, and central to the mechanism claim, the
near-wall ordering is unchanged: \agfno{}'s ring error is
\AgfnoVsUnetRing$\times$ better than the U-Net's and \AgfnoVsCnoRing$\times$
better than CNO's --- the geometry-carrying mechanism, not the Fourier
transform per se, is what the walls reward. Third, CNO --- designed
precisely to inherit operator-style discrete invariance without FFTs ---
does not close the gap on geometry-anchored targets, confirming thatthe failing ingredient of the convolutional family on this benchmark is not
spectral versus spatial smoothing but the absence of an explicit,
per-depth geometry channel.
\fi

% ---------------- Pilot-budget matched baselines (experiments5p) ----------
\ifdefined\PAgfnoRing
\subsection{Pilot-budget matched baselines}
\label{sec:baselines-pilot}

A pilot-budget companion study complements the full-budget comparison:
a uniformly reduced protocol
(\PilotEpochs{} epochs, \PilotNtrain{} training samples, $\PilotRes^2$
grid) applied \emph{identically} to all four architectures, with \fno{} and
\agfno{} retrained under it rather than imported from the headline matrix,
so the four-way ordering is internally valid at this budget. Three seeds
each; the same three metrics and byte-identical splits.

\begin{table}[t]
\centering\small
\caption{Pilot-budget matched baselines (uniformly reduced protocol,
identical for every row; \fno{} and \agfno{} retrained under it).
$n{=}3$ seeds, mean$\pm$std. Ring error and wall fidelity are the
geometry metrics; lower is better in every column.}
\label{tab:baselines-pilot}
\begin{tabular}{lccc}
\toprule
Architecture & \relL{} (ring) & Wall fidelity & \relL{} (global) \\
\midrule
\fno{}          & \PFnoRing$\pm$\PFnoRingStd & \PFnoWall & \PFnoRelL$\pm$\PFnoRelLStd \\
U-Net           & \PUnetRing$\pm$\PUnetRingStd & \PUnetWall & \textbf{\PUnetRelL$\pm$\PUnetRelLStd} \\
CNO             & \PCnoRing$\pm$\PCnoRingStd & \PCnoWall & \PCnoRelL$\pm$\PCnoRelLStd \\
\agfno{} (ours) & \textbf{\PAgfnoRing$\pm$\PAgfnoRingStd} & \textbf{\PAgfnoWall} & \PAgfnoRelL$\pm$\PAgfnoRelLStd \\
\bottomrule
\end{tabular}
\end{table}

Four findings (Table~\ref{tab:baselines-pilot}). First, \agfno{} is the
best architecture on \emph{both} geometry metrics: its ring error is
\PilotUnetRingOverAgf$\times$ lower than the U-Net's,
\PilotFnoRingOverAgf$\times$ lower than the \fno{}'s, and
\PilotCnoRingOverAgf$\times$ lower than CNO's; its wall fidelity is
\PilotAgfWallOverFno$\times$ better than the \fno{}'s and
\PilotAgfWallOverUnet$\times$ better than the U-Net's. Second, an
instructive inversion: the U-Net \emph{beats} the \fno{} on ring error
(\PUnetRing{} vs.\ \PFnoRing) while trailing it globally. Localized
convolution assembles geometry from its receptive field and simply does not
suffer the periodic-global failure mode --- so the ring pathology is not a
generic property of deep networks but is specific to the band-limited
global spectral channel, which is precisely what the mechanism claim
asserts. Third, the mechanistic control holds: CNO, the band-limited global
operator \emph{without} a geometry channel, shows the worst ring error of
the four --- the pathology tracks the global spectral pathway, not the
architecture family. Fourth, the honest disclosure: the U-Net wins global
relative $L^2$ (\PUnetRelL{} vs.\ \PAgfnoRelL{}) at this budget, as a pure
spatial interpolator with global coupling assembled across depth is expected
to on a smooth interior; the paper's claim concerns geometry-anchored
targets, where the U-Net loses on both geometry metrics, and the
full-budget comparison is the decisive test.
\fi


On the literature's piecewise-constant Darcy setting (no boundary penalty,
no ring channel in the loss):
\fno{} \BenchFnoRelL, \agfno{}-frozen \BenchFrozenRelL, \agfno{}
\textbf{\BenchAgfRelL} relative $L^2$ (\BenchAgfVsFnoRed{} vs.\ baseline)---
the same ordering as the obstacle benchmark, obtained with the penalty and
ring-channel confounds removed. For scale only: the published \fno{} number
on the $85^2$ benchmark is \PubFnoDarcy{} and Geo-FNO's is \PubGeoFnoDarcy{}
\citep{li2021fno, mao2023geofno}; these are \emph{not} comparable to ours
(different data generators, resolutions, and splits), and we do not treat
them as baselines.

Deform-\fno{} (learned deformation + \fno{}) reaches \GeofnoRelL{} global /
\GeofnoRing{} ring on obstacle Darcy. The numbers confirm the structural
diagnosis made in Section~\ref{sec:method}: a learned deformation is a
diffeomorphism of the reference frame, but a domain with interior wall
obstacles is \emph{not simply connected}; no deformation can flatten
multiply-connected geometry onto a periodic box. Deform-\fno{} improves the
simply-connected part of the problem (global error \GeofnoRelL{} vs.\
\ExpFnoRelL{} for the undistorted \fno{}) while remaining blind at the
walls (ring \GeofnoRing{}, statistically indistinguishable from the \fno{}'s
\ExpFnoRing{}; wall fidelity \GeofnoWall{} vs.\ \agfno{}'s \ExpAgfWall).
We report this as the honest strongest-effort representative of the
deformation family on this benchmark, not as a refutation of Geo-FNO on its
own simply-connected problems.

\subsection{The forgetting measurement}
\label{sec:probe}

\begin{table}[t]
\centering\small
\caption{Linear-geometry probe: $\Rtwo$ of a closed-form ridge read-out of
the SDF from each \emph{final} block's latent (192 test samples; protocol
with null and sanity controls in the released code).}
\label{tab:probe}
\begin{tabular}{lcccc}
\toprule
Depth & 1 & 2 & 3 & 4 \\
\midrule
\fno{}   & \ProbeFnoDOne & \ProbeFnoDTwo & \ProbeFnoDThree & \ProbeFnoDFour \\
\agfno{} & \ProbeAgfDOne & \ProbeAgfDTwo & \ProbeAgfDThree & \ProbeAgfDFour \\
\bottomrule
\end{tabular}
\end{table}

Table~\ref{tab:probe} and Figure~\ref{fig:probe} are the paper's central
empirical object. In a trained \fno{}, decodable geometry decays
monotonically and steeply with depth ($\Rtwo$ \ProbeRtwoDropFno{} total drop
by depth 4); \agfno{}'s re-injection slows the decay (\ProbeRtwoDropAgf{}
drop; its depth-4 latent retains the geometry access an \fno{} loses by
depth ${\sim}2.5$). We emphasize the honest reading: the mechanism
\emph{mitigates} forgetting at these depths; it does not eliminate it. The
probe comes with a unit-tested null control (noise features
$\Rightarrow \Rtwo \approx 0$) and sanity control (linear embeddings
$\Rightarrow \Rtwo > 0.99$).

\begin{figure}[t]
\centering
\includegraphics[width=0.58\linewidth]{figs/fig_probe.pdf}
\caption{Linear-geometry probe $\Rtwo$ of the final block's latent vs.\
depth. The \fno{}'s boundary decodability collapses; \agfno{}'s
re-injection slows the decay. Null and sanity controls are unit-tested in
the released code.}
\label{fig:probe}
\end{figure}

\begin{figure}[t]
\centering
\includegraphics[width=0.58\linewidth]{figs/fig_depth.pdf}
\caption{Behavioural counterpart: near-boundary error vs.\ depth on
obstacle Darcy. The behavioural curve alone is a weak forgetting test at
these depths---which is precisely why the representational probe
(Figure~\ref{fig:probe}) carries the finding.}
\label{fig:depth}
\end{figure}

\subsection{Depth to 16: collapse, plateau, and the aperture prediction}
\label{sec:deep}

Proposition~\ref{prop:bandlimit}(c) makes a sharp prediction: if forgetting
is the decay of the global channel's effective aperture, then at sufficient
depth the \fno{} should lose geometry \emph{entirely}, while the re-injecting
architecture should reach an equilibrium where re-injection balances decay.
Both arms confirm at depth 16 (Table~\ref{tab:deep}).

\begin{table}[t]
\centering\small
\caption{Extended depth sweep on obstacle Darcy (single seed per cell,
identical budget; probe $\Rtwo$ measured on the final block's latent).
Table~\ref{tab:seeds} replicates the behavioural rows with three seeds.}
\label{tab:deep}
\begin{tabular}{llcccc}
\toprule
 & & depth 4 & depth 6 & depth 8 & depth 16 \\
\midrule
\multirow{2}{*}{\relL{} (global)}
 & \fno{}   & \DFourFnoRelL & \DSixFnoRelL & \DEightFnoRelL & \DSixteenFnoRelL \\
 & \agfno{} & \DFourAgfnoRelL & \DSixAgfnoRelL & \DEightAgfnoRelL & \textbf{\DSixteenAgfnoRelL} \\
\midrule
\multirow{2}{*}{\relL{} (ring)}
 & \fno{}   & \DFourFnoRing & \DSixFnoRing & \DEightFnoRing & \DSixteenFnoRing \\
 & \agfno{} & \DFourAgfnoRing & \DSixAgfnoRing & \DEightAgfnoRing & \textbf{\DSixteenAgfnoRing} \\
\midrule
\multirow{2}{*}{probe $\Rtwo$}
 & \fno{}   & \ProbeExtFnoDFour & \ProbeExtFnoDSix & \ProbeExtFnoDEight & \textbf{\ProbeExtFnoDSixteen} \\
 & \agfno{} & \ProbeExtAgfDFour & \ProbeExtAgfDSix & \ProbeExtAgfDEight & \ProbeExtAgfDSixteen \\
\bottomrule
\end{tabular}
\end{table}

\textbf{The \fno{} collapses.} At depth 16 its global error is
\DSixteenFnoRelL{}---\emph{worse than predicting the training mean}---and its
final-block probe $\Rtwo$ is \ProbeExtFnoDSixteen{}, i.e.\ the latent carries
no linearly decodable geometry at all. Behavioural collapse and
representational collapse coincide, which is exactly what the aperture
mechanism requires: the network's only long-range channel has nothing left
to say about where the walls are.

\textbf{\agfno{} plateaus---and improves.} Its probe $\Rtwo$ holds within
\ProbeExtAgfPlateauSpread{} across depths 4--16 (the predicted plateau), and
its \emph{behavioural} error is flat to improving with depth (ring error
\DFourAgfnoRing{} $\to$ \DSixteenAgfnoRing{}): with geometry continuously
re-supplied, extra depth buys solution capacity instead of forgetting.
The honest reading stands: the mechanism converts depth from a liability
into a resource for wall-anchored problems; it does not make the probe's
$\Rtwo$ recover to shallow-network levels, and the plateau is equilibrium,
not elimination.

\begin{table}[t]
\centering\small
\caption{Seed replication of the depth sweep (mean $\pm$ std over three
seeds per cell; identical budget per run).}
\label{tab:seeds}
\begin{tabular}{llccc}
\toprule
 & & depth 4 & depth 8 & depth 16 \\
\midrule
\multirow{2}{*}{\relL{} (global)}
 & \fno{}   & \SweepDFourFnoRelL\,$\pm$\,\SweepDFourFnoRelLStd & \SweepDEightFnoRelL\,$\pm$\,\SweepDEightFnoRelLStd & \SweepDSixteenFnoRelL\,$\pm$\,\SweepDSixteenFnoRelLStd \\
 & \agfno{} & \textbf{\SweepDFourAgfnoRelL\,$\pm$\,\SweepDFourAgfnoRelLStd} & \SweepDEightAgfnoRelL\,$\pm$\,\SweepDEightAgfnoRelLStd & \SweepDSixteenAgfnoRelL\,$\pm$\,\SweepDSixteenAgfnoRelLStd \\
\midrule
\multirow{2}{*}{\relL{} (ring)}
 & \fno{}   & \SweepDFourFnoRing\,$\pm$\,\SweepDFourFnoRingStd & \SweepDEightFnoRing\,$\pm$\,\SweepDEightFnoRingStd & \SweepDSixteenFnoRing\,$\pm$\,\SweepDSixteenFnoRingStd \\
 & \agfno{} & \SweepDFourAgfnoRing\,$\pm$\,\SweepDFourAgfnoRingStd & \SweepDEightAgfnoRing\,$\pm$\,\SweepDEightAgfnoRingStd & \textbf{\SweepDSixteenAgfnoRing\,$\pm$\,\SweepDSixteenAgfnoRingStd} \\
\bottomrule
\end{tabular}
\end{table}

\textbf{The collapse replicates.} Re-running the sweep with three seeds per
cell (Table~\ref{tab:seeds}) leaves every qualitative conclusion intact and
makes the collapse essentially deterministic: at depth 16 the \fno{} sits at
\SweepDSixteenFnoRelL{}\,$\pm$\,\SweepDSixteenFnoRelLStd{} (worst seed
\SweepGOneFnoWorst{}), a $\SweepFnoAgfSixteenRatio\times$ error gap with
\emph{zero} overlap across seeds, while \agfno{} varies by at most
$\SweepAgfSixteenRelVar$ of its mean. Forgetting at depth is not a tail
event; it is what the architecture does.

\subsection{Frequency-resolved analysis}
\label{sec:analysis}

Testing P1--P2 directly on trained checkpoints
(\texttt{agfno/analysis.py}, results regenerated from the released weights).
The premise: the SDF's spectrum carries \FreqSdfLowEnergy{} of its energy
below the architectural cut, \FreqSdfHighEnergy{} above---the direct tail is
real but tiny, exactly as the $|k|^{-2}$ argument predicts.

\textbf{P2 (spectral signature) confirms.} Band-resolved ring error:
\fno{}'s misfit piles up in the truncated band (high-band ring error
\FreqFnoHighRing{}, high-over-low ratio \FreqFnoHighOverLow{}), and \agfno{}
reduces precisely that band---\FreqAgfHighRingRed{} lower high-band ring
error, with the high-over-low ratio halved to \FreqAgfHighOverLow{}. No
other band shows a comparable relative reduction; the mechanism acts where
the proposition says it must.

\textbf{P1 (propagation) is directionally consistent.} A wall-localized
perturbation of the input propagates with a more wall-confined influence
through the \fno{} (ring-over-far influence ratio \FreqPoneFnoRingOverFar{})
than through \agfno{} (\FreqPoneAgfRingOverFar{}), consistent with the
sinc-kernel confinement picture; the full spatial decay profiles are in the
released artifacts.

\subsection{Three dimensions: the collapse--plateau replication in 3-D}
\label{sec:3d}

Every result above lives in 2-D. The mechanism claims, however, are about
the operator's band structure and the Markovian decay of geometric
information -- neither of which knows about dimension. We therefore lifted
the full stack to 3-D obstacle Darcy: the unit cube with random heterogeneous
permeability and 1--2 analytic obstacles per sample, now spheres \emph{and}
solid tori (the torus keeps the domain multiply-connected in 3-D, where the
deformation-based escape is structurally unavailable), exact Dirichlet
walls, ground truth from a float64 PCG solver (every stored field satisfies
the discrete PDE to $<$10$^{-6}$ relative residual), and zero-gate capacity
control carried over unchanged. The 3-D operator uses width 32 with an
8$^3$ mode cut (the 3-D spectral weights and FFTs dominate cost); everything
else -- protocol, seeds, budget -- matches the 2-D sweep.

\begin{table}[t]
\centering\small
\caption{3-D replication (res 32, mean $\pm$ std over \ThreeDSeeds{}
seeds). The collapse of vanilla \fno{} at depth 16 and the depth-stability
of \agfno{} reproduce in 3-D; the probe separates the two phenomena
representationally as well as behaviourally. SR = zero-shot super-resolution
to res 48.}
\label{tab:3d}
\begin{tabular}{llcccc}
\toprule
 & & global & near-wall & wall viol. & probe $\Rtwo$ \\
\midrule
\multirow{2}{*}{depth 4}
 & \fno{}   & \ThreeDDFourFnoRelL\,$\pm$\,\ThreeDDFourFnoRelLStd
   & \ThreeDDFourFnoRing\,$\pm$\,\ThreeDDFourFnoRingStd
   & \ThreeDDFourFnoWall & \ThreeDDFourFnoProbe\,$\pm$\,\ThreeDDFourFnoProbeStd \\
 & \agfno{} & \ThreeDDFourAgfnoRelL\,$\pm$\,\ThreeDDFourAgfnoRelLStd
   & \ThreeDDFourAgfnoRing\,$\pm$\,\ThreeDDFourAgfnoRingStd
   & \ThreeDDFourAgfnoWall & \ThreeDDFourAgfnoProbe\,$\pm$\,\ThreeDDFourAgfnoProbeStd \\
\midrule
\multirow{2}{*}{depth 16}
 & \fno{}   & \ThreeDDSixteenFnoRelL\,$\pm$\,\ThreeDDSixteenFnoRelLStd
   & \ThreeDDSixteenFnoRing\,$\pm$\,\ThreeDDSixteenFnoRingStd
   & \ThreeDDSixteenFnoWall & \ThreeDDSixteenFnoProbe\,$\pm$\,\ThreeDDSixteenFnoProbeStd \\
 & \agfno{} & \ThreeDDSixteenAgfnoRelL\,$\pm$\,\ThreeDDSixteenAgfnoRelLStd
   & \ThreeDDSixteenAgfnoRing\,$\pm$\,\ThreeDDSixteenAgfnoRingStd
   & \ThreeDDSixteenAgfnoWall & \ThreeDDSixteenAgfnoProbe\,$\pm$\,\ThreeDDSixteenAgfnoProbeStd \\
\bottomrule
\end{tabular}
\end{table}

\begin{figure}[t]
\centering
\includegraphics[width=0.58\linewidth]{figs/fig_3d.pdf}
\caption{Near-boundary error vs.\ depth in 3-D obstacle Darcy (mean $\pm$ std
over three seeds). Vanilla \fno{} collapses by depth 16; \agfno{} is
flat -- the 2-D forgetting curve, reproduced with spheres and tori.}
\label{fig:3d}
\end{figure}

Table~\ref{tab:3d} reports the 3-D matrix at depths 4 and 16 (three seeds
each). Three observations. \textbf{(i) The collapse is not a 2-D artifact:}
vanilla \fno{} degrades from \ThreeDDFourFnoRelL{} (depth 4) to
\ThreeDDSixteenFnoRelL{} (depth 16) globally, with wall violations rising to
\ThreeDDSixteenFnoWall{} -- the same total-collapse signature as the 2-D
sweep, now with interior toroidal walls the network has never been able to
flatten. \textbf{(ii) The plateau reproduces:} \agfno{} at depth 16 holds
\ThreeDDSixteenAgfnoRelL{} globally and \ThreeDDSixteenAgfnoRing{} at the
walls (vs.\ \ThreeDDFourAgfnoRing{} at depth 4), and its zero-shot
super-resolution to res 48 (\ThreeDDSixteenAgfnoSR{} at depth 16) confirms
the transfer survives the third dimension. \textbf{(iii) The representation
collapses too:} the final-block probe reads $\Rtwo =
\ThreeDDSixteenFnoProbe$ from the \fno{} latent -- near-zero decodable
geometry -- against \ThreeDDSixteenAgfnoProbe{} for \agfno{}, mirroring the
2-D collapse/plateau pair. All pre-registered stability gates pass
(\ThreeDGatesPass{}). Geometric forgetting in mode-truncated operators, and
its restoration by zero-gated geometry injection, is dimension-independent.

% ---------------- Third architecture: AFNO sweep (experiments6) ------------
\ifdefined\AfnoCollapseRatio
\subsection{Architecture generality: the \emph{afno} family member}
\label{sec:afno}

The claim under test is about a \emph{family} -- band-limited periodic
global channels -- not about one implementation. The \afno{} (Adaptive
Fourier Neural Operator) \citep{guibas2022afno}, as deployed at scale in
FourCastNet \citep{pathak2022fourcastnet}, replaces the \fno{}'s hard mode
truncation and post-norm residual design with soft spectrum shrinkage and a
pre-norm, double-skip Transformer-style block, while keeping the same
mode-restricted global spectral pathway and the same parameter budget
(\AfnoParams{} vs.\ \FnoParams{} -- ratio \AfnoParamsRatio). We therefore
re-run the full multi-seed depth protocol of Section~\ref{sec:collapse}
unchanged -- depths 4/8/16, three seeds, byte-identical data and budget --
on the \afno{} and on \emph{agf-afno}, the exact transplant of the
zero-gated geometry mechanism onto the \afno{} block (spectral modulation
plus pointwise re-injection; both gates zero-initialized, so at
initialization the transplant is the plain \afno{}).

\begin{table}[t]
\centering\small
\caption{Third architecture (\afno{} family member): the depth sweep of
Section~\ref{sec:collapse} re-run on the \afno{} and the mechanism
transplant \emph{agf-afno} (mean $\pm$ std over 3 seeds; identical data,
budget, and protocol). The collapse--plateau pair is architecture-generic,
and the transplant rescues it.}
\label{tab:afno}
\begin{tabular}{llcccc}
\toprule
 & & global & near-wall & wall viol. & probe $\Rtwo$ \\
\midrule
\multirow{2}{*}{depth 4}
 & \afno{}     & \AfnoFourRelL\,$\pm$\,\AfnoFourRelLStd
   & \AfnoFourRing\,$\pm$\,\AfnoFourRingStd
   & \AfnoFourWall & \AfnoFourProbe \\
 & \emph{agf-afno} & \AgfafnoFourRelL\,$\pm$\,\AgfafnoFourRelLStd
   & \AgfafnoFourRing\,$\pm$\,\AgfafnoFourRingStd
   & \AgfafnoFourWall & \AgfafnoFourProbe \\
\midrule
\multirow{2}{*}{depth 16}
 & \afno{}     & \AfnoSixteenRelL\,$\pm$\,\AfnoSixteenRelLStd
   & \AfnoSixteenRing\,$\pm$\,\AfnoSixteenRingStd
   & \AfnoSixteenWall & \AfnoSixteenProbe \\
 & \emph{agf-afno} & \AgfafnoSixteenRelL\,$\pm$\,\AgfafnoSixteenRelLStd
   & \AgfafnoSixteenRing\,$\pm$\,\AgfafnoSixteenRingStd
   & \AgfafnoSixteenWall & \AgfafnoSixteenProbe \\
\bottomrule
\end{tabular}
\end{table}

Four pre-registered gates, all passing (\ExpSixGatesPass).
\textbf{(A1) The \afno{} collapses at depth 16}: mean global error
\AfnoSixteenRelL{} with the worst seed no better than \AfnoSixteenWorst{}
-- the total-collapse signature, on a mixer that never hard-truncates a
mode and normalizes like a Transformer. The failure is therefore not an
artifact of the \fno{}'s specific normalization or skip design: band
limitation of the periodic global channel is sufficient.
\textbf{(A2) Depth is a resource once geometry is re-injected}:
\emph{agf-afno} at depth 16 (\AgfafnoSixteenRelL) is its own best depth
cell. \textbf{(A3) The transplant rescues the collapse}:
\AfnoCollapseRatio$\times$ at depth 16 globally
(\AfnoRingRatio$\times$ at the walls). \textbf{(A4) The mechanism, not
just the symptom}: the linear geometry probe's final-block $\Rtwo$ drops by
\AfnoProbeDrop{} across the sweep for the plain \afno{} but only
\AgfafnoProbeDrop{} for the transplant -- the representation-level
separation reproduces on a third architecture with a different
normalization, skip, and spectrum treatment. Together with the depth-16
results of Section~\ref{sec:collapse} this covers the two canonical
designs of the spectral-operator literature; geometric forgetting, and its
zero-gated restoration, are properties of the band-limited global channel
class.
\fi

\section{Reproducibility}
\label{sec:repro}

Every experiment is a self-contained kernel with the full package embedded;
each reruns from a clean container with zero internet access and produces the
checkpoint set, metrics JSONs, and figures consumed by this paper. Data are
byte-deterministic: local regeneration reproduces the remote splits exactly,
and local re-evaluation of the pulled checkpoints matches the remote metrics
to four decimals. Public artifacts: training/checkpoint repository
(\texttt{huggingface.co/Sejibeji/agfno-darcy}), and one-click reruns of the
headline, controlled-suite, benchmark, extended, and 3-D replication
experiments (\texttt{kaggle.com/code/sehajrsingh/agfno-darcy-*}).

\section{Limitations}
\label{sec:limitations}

(i) Two PDE families; the geometry suite spans polygonal obstacle fields in
2-D and spheres and solid tori in 3-D, but airfoil-class geometries and
time-varying boundary motion remain untested. (ii) Published numbers are
quoted for scale with explicit non-comparability caveats; external
validation comes from U-Net and CNO reimplemented at matched settings on
our splits (Section~\ref{sec:baselines-pilot}\ifdefined\UNetRelL{}
and Section~\ref{sec:baselines}\fi;
the pilot establishes the four-way ordering at a uniformly reduced
protocol, the full-budget comparison is the decisive one) --- faithful
re-implementations, not the authors' original code, and no external
\emph{operator} baseline beyond the convolutional and deformation families
is included. (iii) The probe measures \emph{linear}
decodability---a lower bound on retained geometric information; the
mitigation-vs-elimination distinction at depth 16 relies on it.
(iv) Obstacles act as fixed-temperature (PDE2) or sealed (Darcy)
inclusions; dynamic boundary motion is out of scope. (v) Single-family
velocity field (smooth random stream function) for PDE2.
(vi) The $\Delta\rho$ diagnostic is validated on the two PDE families it
was constructed to explain; it demonstrates functional form, not a
validated selection rule, until tested on further datasets
(Section~\ref{sec:diagnostic}).

\section{Conclusion}

Geometric forgetting in spectral operators is not mysterious; it is the
band-limitation of their only global channel, stated in two lines and
measurable in one. A zero-gated SDF modulation restores the truncated band
by spectral convolution, converts the shift-invariant kernel into a
wall-localized one, and delivers a controlled, seed-replicated,
cross-PDE-validated error reduction whose attribution survives capacity,
component, and loss ablations. The deeper contribution is methodological:
the hypothesis that operators forget geometry has circulated as folklore;
here it acquires a number, a control, and a falsifiable plateau prediction.

\bibliographystyle{plainnat}
\begin{thebibliography}{25}

\bibitem[Alain and Bengio(2017)]{alain2017probing}
G.~Alain, Y.~Bengio.
\newblock Understanding intermediate layers using linear classifier probes.
\newblock arXiv:1610.01644, 2017.

\bibitem[Basri et~al.(2020)]{basri2020frequency}
R.~Basri, M.~Galun, A.~Geifman, D.~Jacobs, Y.~Kasten, S.~Kritchman.
\newblock The convergence rate of neural networks for learned functions of
different frequencies.
\newblock \emph{NeurIPS}, 2020.

\bibitem[Bi et~al.(2023)]{bi2023graphcast}
K.~Bi, L.~Xie, H.~Zhang, X.~Chen, et~al.
\newblock Accurate medium-range global weather forecasting with 3D neural
networks.
\newblock \emph{Nature}, 619, 2023.

\bibitem[Brandstetter et~al.(2022)]{brandstetter2022message}
J.~Brandstetter, D.~Worrall, M.~Welling.
\newblock Message Passing Neural PDE Solvers.
\newblock \emph{ICLR}, 2022.

\bibitem[Cao(2021)]{cao2021galerkin}
S.~Cao.
\newblock Choose a Transformer: Fourier or Galerkin.
\newblock \emph{NeurIPS}, 2021.

\bibitem[Hao et~al.(2023)]{hao2023gnot}
Z.~Hao, S.~Liu, Y.~Zhang, C.~Ying, Y.~Feng, H.~Su, J.~Zhu.
\newblock GNOT: A General Neural Operator Transformer for Operator Learning.
\newblock \emph{ICML}, 2023.

\bibitem[Guibas et~al.(2022)]{guibas2022afno}
J.~Guibas, M.~Mardani, Z.~Li, A.~Tao, A.~Anandkumar, B.~Catanzaro.
\newblock Adaptive Fourier Neural Operators: Efficient Token Mixers for
Transformers.
\newblock \emph{ICLR}, 2022.

\bibitem[Kovachki et~al.(2023)]{kovachki2023nocal}
N.~Kovachki, Z.~Li, B.~Liu, K.~Azizzadenesheli, et~al.
\newblock Neural Operator: Learning Maps Between Function Spaces.
\newblock \emph{JMLR}, 24(89), 2023.

\bibitem[Lam et~al.(2023)]{lam2023learning}
R.~Lam, A.~Sanchez-Gonzalez, M.~Willson, P.~Wirnsberger, et~al.
\newblock Learning skillful medium-range global weather forecasting.
\newblock \emph{Science}, 382, 2023.

\bibitem[Li et~al.(2021)]{li2021fno}
Z.~Li, N.~Kovachki, K.~Azizzadenesheli, B.~Liu, K.~Bhattacharya, A.~Stuart,
A.~Anandkumar.
\newblock Fourier Neural Operator for Parametric Partial Differential
Equations.
\newblock \emph{ICLR}, 2021. arXiv:2010.08895.

\bibitem[Li et~al.(2023)]{li2021grafno}
Z.~Li, D.~Zheng, N.~Kovachki, D.~Jin, H.~Chen, B.~Liu, K.~Azizzadenesheli,
A.~Anandkumar.
\newblock Physics-Informed Neural Operator for Learning Partial Differential
Equations.
\newblock \emph{ACM MSML}, 2024. (Graph-NO line.)

\bibitem[Lu et~al.(2021)]{lu2021deeponet}
L.~Lu, P.~Jin, G.~E.~Karniadakis.
\newblock DeepONet: Learning nonlinear operators via DeepONet based on the
universal approximation theorem of operators.
\newblock \emph{Nature Machine Intelligence}, 3, 2021.

\bibitem[Mao et~al.(2023)]{mao2023geofno}
Z.~Mao, Z.~Li, L.~Lu, T.~K.~Ghosh, G.~E.~Karniadakis.
\newblock Fourier Neural Operator with Learned Deformations for PDEs on
General Geometric Domains.
\newblock \emph{JMLR}, 24(388), 2023.

\bibitem[Park et~al.(2019)]{park2019deepsdf}
J.~J.~Park, P.~Florence, J.~Straub, R.~Newcombe, S.~Lovegrove.
\newblock DeepSDF: Learning Continuous Signed Distance Functions for Shape
Representation.
\newblock \emph{CVPR}, 2019.

\bibitem[Pathak et~al.(2022)]{pathak2022fourcastnet}
J.~Pathak, S.~Subramanian, P.~Harrington, S.~Raja, A.~Chattopadhyay, et~al.
\newblock FourCastNet: A Global Data-driven High-resolution Weather Model
using Adaptive Fourier Neural Operators.
\newblock arXiv:2202.11214, 2022.

\bibitem[Perez et~al.(2018)]{perez2018film}
E.~Perez, F.~Strub, H.~de~Vries, V.~Dumoulin, A.~Courville.
\newblock FiLM: Visual Reasoning with a General Conditioning Layer.
\newblock \emph{AAAI}, 2018.

\bibitem[Pfaff et~al.(2021)]{pfaff2021learning}
T.~Pfaff, M.~Fortunato, A.~Sanchez-Gonzalez, P.~Battaglia.
\newblock Learning Mesh-Based Simulation with Graph Networks.
\newblock \emph{ICLR}, 2021.

\bibitem[Raoni{\'c} et~al.(2024)]{raonic2024cno}
B.~Raoni{\'c}, R.~Molinaro, T.~Rohner, K.~Kashinath, M.~Keller.
\newblock Convolutional Neural Operators for robust and accurate learning of
PDEs.
\newblock \emph{NeurIPS}, 2024.

\bibitem[Rahaman et~al.(2019{a})]{rahaman2019spectral}
N.~Rahaman, A.~Baratin, D.~Arpit, F.~Draxler, M.~Lin, F.~Hamprecht,
Y.~Bengio, A.~Courville.
\newblock On the Spectral Bias of Neural Networks.
\newblock \emph{ICML}, 2019.

\bibitem[Rahaman et~al.(2019{b})]{rahaman2019cnn}
N.~Rahaman, D.~Arpit, A.~Baratin, F.~Draxler, et~al.
\newblock On the Spectral Bias of Convolutional Neural Networks.
\newblock \emph{ICML}, 2019.

\bibitem[Ronneberger et~al.(2015)]{ronneberger2015unet}
O.~Ronneberger, P.~Fischer, T.~Brox.
\newblock U-Net: Convolutional Networks for Biomedical Image Segmentation.
\newblock \emph{MICCAI}, 2015.

\bibitem[Rusch and Bronstein(2023)]{rusch2023survey}
T.~K.~Rusch, M.~M.~Bronstein.
\newblock A Survey on Over-smoothing in Graph Neural Networks.
\newblock arXiv:2303.10993, 2023.

\bibitem[Singh(2026)]{singh2026agfno}
S.~R.~Singh.
\newblock AGF-NO: code, checkpoints, and one-click GPU reruns.
\newblock \url{https://huggingface.co/Sejibeji/agfno-darcy}, 2026.

\bibitem[Takamoto et~al.(2022)]{takamoto2022pdebench}
M.~Takamoto, T.~Praditia, R.~Leiteritz, D.~MacKinlay, F.~Alesiani,
M.~Pfl{\"u}ger, N.~Thuerey.
\newblock PDEBench: An Extensive Benchmark for Scientific Machine Learning.
\newblock \emph{NeurIPS Datasets and Benchmarks}, 2022.

\bibitem[Wang et~al.(2022)]{wang2022freq}
Z.-Q.~J.~Wang, X.~Ma, Z.~Xu.
\newblock Frequency Principle: Fourier Analysis Sheds Light on Deep Neural
Networks.
\newblock \emph{Communications in Computational Physics}, 28(5), 2022.

\bibitem[Xia and Aviles-Rivero(2026)]{xia2026forgetting}
Y.~Xia, A.~I.~Aviles-Rivero.
\newblock Do Neural Operators Forget Geometry? The Forgetting Hypothesis in
Deep Operator Learning.
\newblock \emph{ICML}, 2026. arXiv:2605.05862. (Concurrent work; see
Section~\ref{sec:related}.)

\bibitem[Yin et~al.(2023)]{yin2023gino}
Y.~Yin, M.~Kirchdoerfer, S.~M.~Kaplan, et~al.
\newblock Geometry-Informed Neural Operator for Large-Scale 3D PDEs.
\newblock \emph{NeurIPS}, 2023.

\end{thebibliography}

\end{document}
